%%%%%%%%%%  scan idx 56  |  folio 45  %%%%%%%%%%

\textbf{1.15.} En lo que resta de este párrafo explicaremos brevemente
cómo modificar la teoría anterior para permitir extensiones del cuerpo
residual.

Sean $\Lambda$ y $K$ anillos locales noetherianos completos, siendo
$K$ una $\Lambda$-álgebra. Supongamos que:

(a) el morfismo natural $\Lambda\longrightarrow K$ es un
homomorfismo local;

(b) el cuerpo residual $K'$ de $K$ es una extensión de tipo finito del
cuerpo residual $k$ de $\Lambda$.

El caso más interesante es $K=K'$, siendo $K$ una extensión finita
inseparable de $k$.

Denotemos por $\widehat C_{\Lambda,K}$ la categoría cuyos objetos son
anillos locales noetherianos completos $R$ provistos de una estructura
de $\Lambda$-álgebra y de un epimorfismo de $\Lambda$-álgebras
$s:R\longrightarrow K$. Sea $C_{\Lambda,K}$ su subcategoría plena
constituida por los pares $(R,s)$ para los cuales $\ker(s)$ tiene
longitud finita como $R$-módulo. La categoría
$\widehat C_{\Lambda,K}$ se identifica con una subcategoría plena de
$\operatorname{Pro}(C_{\Lambda,K})$. Si $K$ es artiniano, todo objeto
de $C_{\Lambda,K}$ es artiniano.

Para un espacio vectorial $V$ de dimensión finita sobre $K'$, denotemos
por $D_V(K)$ el álgebra de números duales sobre $K$ definida por $V$.
Sus elementos son de la forma $x+v\varepsilon$, donde $x\in K$,
$v\in V$ y $\varepsilon^2=0$, con multiplicación
\[
 (x+v\varepsilon)(x'+v'\varepsilon)
 =xx'+(xv'+x'v)\varepsilon.
\]
El morfismo
\[
 s:D_V(K)\longrightarrow K,
 \qquad s(x+v\varepsilon)=x,
\]
hace de $D_V(K)$ un objeto de $C_{\Lambda,K}$. El papel que antes
desempeñaba $k[\varepsilon]$ lo desempeña ahora $D_{K'}(K)$.

Como en 1.8, todo grupoide cofibrado $\underline A$ sobre
$C_{\Lambda,K}$ se extiende naturalmente a una categoría cofibrada
$\widehat{\underline A}$ sobre $\widehat C_{\Lambda,K}$. Un morfismo
de $\underline A$ o $\widehat{\underline A}$ se llama
\underline{sobreyectivo} cuando el homomorfismo de álgebras
correspondiente en $C_{\Lambda,K}$ o $\widehat C_{\Lambda,K}$ es
sobreyectivo.

%%%%%%%%%%  scan idx 57  |  folio 46  %%%%%%%%%%

\underline{Suponemos que $\underline A(K)=\{e\}$.}

\underline{Definición 1.16. Un grupoide cofibrado $\underline A$ sobre
$C_{\Lambda,K}$ se llama semihomogéneo si se cumplen las propiedades
siguientes:}

(I) \underline{Todo diagrama de flechas continuas}
\[
\begin{tikzcd}[row sep=large,column sep=large]
a'' \arrow[r,dashed] \arrow[d,dashed]
 & a' \arrow[d] \\
a_1 \arrow[r] & a
\end{tikzcd}
\]
\underline{en $\underline A$, donde $a'\longrightarrow a$ es
sobreyectivo, puede completarse en un diagrama conmutativo añadiendo
las flechas punteadas.}

(II) \underline{Para todo $R$, el morfismo natural}
\[
 [\underline A(R\times_KD_{K'}(K))]
 \longrightarrow
 [\underline A(R)]\times[\underline A(D_{K'}(K))]
\]
\underline{es biyectivo.}

Si $\underline A$ es semihomogéneo, también lo es $[\underline A]$.

\textbf{1.17.} Restrinjamos primero $\underline A$ a la subcategoría
plena de $C_{\Lambda,K}$ cuyos objetos son los $D_V(K)$. Pongamos
\[
 \Omega=\Omega^1_{K/\Lambda}\otimes_KK',
 \qquad d:K\longrightarrow\Omega.
\]
Un morfismo $f:D_V(K)\longrightarrow D_W(K)$ puede escribirse
\[
 f(x+v\varepsilon)
 =x+\bigl(u(v)+D(dx)\bigr)\varepsilon,
\]
con $u\in\operatorname{Hom}_{K'}(V,W)$ y
$D\in\operatorname{Hom}_{K'}(\Omega,W)$. Identificando $f$ con
$(u,D)$, la composición viene dada por
\[
 (u,D)\circ(u',D')=(uu',D+uD').
\]
La condición (II) implica:

(II$'$) el funtor de los espacios vectoriales sobre $K'$ a los
conjuntos
\[
 V\longmapsto[\underline A(D_V(K))]
\]
conmuta con los productos.

